\documentclass[12pt]{article}
%% BUGS:
%% - The \begin{problem} ... \end{problem} is only used in exams??

\usepackage{url, graphicx, epstopdf}

% page layout
\usepackage[letterpaper, textheight=9.5in, textwidth=5.5in]{geometry}
\usepackage{setspace}
\setstretch{1.08}
\pagestyle{plain}
\sloppy\sloppypar\raggedbottom\frenchspacing\thispagestyle{empty}

% problem set header
\newcommand{\psname}{Problem Set}
\newcommand{\startps}[1]{\section*{Computational Physics --- \psname~{#1}}}

% problem formatting
\usepackage{amsthm}
\theoremstyle{definition}
\newcommand{\problemname}{Problem}
\newtheorem{problem}{\problemname}
\newcommand{\probref}[1]{\problemname~\ref{#1}}
\theoremstyle{remark}
\newcommand{\notename}{Note}
\newtheorem{note}{\notename}[problem]
\newcommand{\noteref}[1]{\textit{\notename}~\ref{#1}}
\newcounter{subpart}[problem]
\newcommand{\subpart}{\par\addvspace{1ex}\refstepcounter{subpart}\textsl{(\alph{subpart})}~}

% words
\newcommand{\foreign}[1]{\textsl{#1}}
\newcommand{\vs}{\foreign{vs}}

% code
\newcommand{\code}[1]{\texttt{\detokenize{#1}}}
\usepackage{xcolor}
% Define custom colors for code highlighting
\definecolor{codegreen}{rgb}{0,0.6,0}
\definecolor{codegray}{rgb}{0.5,0.5,0.5}
\definecolor{codepurple}{rgb}{0.58,0,0.82}
\definecolor{backcolour}{rgb}{0.99,0.99,0.99}
\usepackage{listings}
\lstdefinestyle{pythonstyle}{
    backgroundcolor=\color{backcolour},   
    commentstyle=\color{codegreen},
    keywordstyle=\color{magenta},
    numberstyle=\tiny\color{codegray},
    stringstyle=\color{codepurple},
    tabsize=4
}
\lstset{
    basicstyle=\small\ttfamily,
    columns=fixed,
    breaklines=false,
    keepspaces=true,
    showspaces=false,
    showstringspaces=false,
    showtabs=false,
    showlines=*,
    style=pythonstyle,
    language=Python,
    literate={~}{{\url{~}}}1
}

% math
\usepackage{amsmath}
\newcommand{\dd}{\mathrm{d}}
\newcommand{\e}{\mathrm{e}}

% primary units
\newcommand{\rad}{\mathrm{rad}}
\newcommand{\kg}{\mathrm{kg}}
\newcommand{\m}{\mathrm{m}}
\newcommand{\s}{\mathrm{s}}

% secondary units
\renewcommand{\deg}{\mathrm{deg}}
\newcommand{\km}{\mathrm{km}}
\newcommand{\cm}{\mathrm{cm}}
\newcommand{\mm}{\mathrm{mm}}
\newcommand{\ft}{\mathrm{ft}}
\newcommand{\mi}{\mathrm{mi}}
\newcommand{\AU}{\mathrm{AU}}
\newcommand{\ns}{\mathrm{ns}}
\newcommand{\h}{\mathrm{h}}
\newcommand{\yr}{\mathrm{yr}}
\newcommand{\N}{\mathrm{N}}
\newcommand{\J}{\mathrm{J}}
\newcommand{\eV}{\mathrm{eV}}
\newcommand{\W}{\mathrm{W}}
\newcommand{\Pa}{\mathrm{Pa}}

% derived units
\newcommand{\mps}{\m\,\s^{-1}}
\newcommand{\mph}{\mi\,\h^{-1}}
\newcommand{\mpss}{\m\,\s^{-2}}

% random units
\newcommand{\au}{\mathrm{a.u.}}


\begin{document}

\startps{4}
Due Friday 2026 October 2 at 18:00 on Brightspace.
Be sure to explicitly give credit where it is due, on all problems.
Also credit buddies, web pages, and LLMs (as per course policy).
Failure to give explicit credits will result in points deducted.

\begin{problem}
  Grab the data at \hogg{[THIS LINK]}; it is a pickle file.
  Read in the pickle file with something like this:
  \begin{lstlisting}
    xs, ys = pickle.load(fn)\end{lstlisting}
  \subpart Find the value of $x$ in \code{xs} that corresponds to the maximum value of $y$ in \code{ys}.
  \subpart Now fit a parabola to $\ln(y)$ as a function of $x$ using the three points centered on the maximum point you found in the previous part.
  You could use high school math methods to do this, or you might use \code{numpy.polynomial} code.
  Find the value of $x$ corresponding to the maximum of that parabola.
  \subpart Plot the input data as black dots on a properly formatted and labeled plot.
  Plot the best-fit parabola as a thin curve.
  Plot a vertical line at the maximum you obtained in the previous part.
  Is everything cool?
\end{problem}

\begin{note}
  In the last part (the plotting), you might have to think about the fact that you fit the polynomial in $\ln(y)$ but are plotting $x$ \foreign{vs} $y$.
\end{note}

\begin{note}
  When you plot the parabola, you want to make a dense set of $x$ values so the parabola looks like a nice smooth curve.
\end{note}

\begin{note}
  If you search GitHub(tm) you can find Hogg's code that made the fake data.
  Is your answer correct?
  If not, why not?
  These questions are not for credit!
\end{note}

\begin{problem}
  Newman, \textit{Computational Physics, 2ed}, Exercise 5.12.
\end{problem}

\begin{problem}
  foo
\end{problem}

\begin{problem}
  foo
\end{problem}

\end{document}
